\documentclass[11pt]{article}
\usepackage{amsmath,amssymb,geometry,enumitem}
\geometry{margin=1in}
\title{BindingAndCatalysis.jl: Regime Canonicalization, Provenance, and Singular-Aware Mapping Refactor}
\author{Internal Technical Report}
\date{\today}
\begin{document}
\maketitle

\section{Problem Statement}
The prior regime representation stored dense per-regime constraints and mappings as primary ownership. This duplicated three levels of structure: (i) repeated $x$-space dominance boundaries, (ii) closely related mapping objects across neighboring regimes, and (iii) transformed $qK$-space inequalities. Most importantly, geometric provenance was not first-class: one could not reliably query which $x$-space boundary atom, transported through which mapping object, generated a specific $qK$ interface.

\section{Refactor Goals}
We implemented a structural model preserving the causal chain
\[
\text{$x$-halfspace atom}\to\text{mapping object}\to\text{$qK$ atom},
\]
while keeping existing materialized APIs compatible.

\section{New Architecture}
\subsection{Layer A: Canonical $x$-space atoms}
A new global pool stores unique undirected halfspace atoms keyed by
\((i,\min(j_1,j_2),\max(j_1,j_2))\).
Each atom stores one canonical direction and offset, while regime-local orientation is stored separately as an \texttt{Int8} sign.

\subsection{Layer B: Mapping-object pool}
Each vertex now references a canonical mapping ID. Mapping objects carry nullity, $H$, $H_0$ (when defined), and $M_0$ plus singular-normal metadata for nullity-1 interpretability.

\subsection{Layer C: Canonical $qK$ atoms with provenance}
A $qK$ atom is pooled by the key
\((\texttt{x\_halfspace\_id},\texttt{orientation},\texttt{mapping\_id})\).
Hence every inequality/hyperplane in projected space carries explicit provenance and can be traced back without recomputing dense matrices.

\subsection{Layer D: Lightweight regime ownership}
\texttt{Vertex} stores references:
\begin{itemize}[leftmargin=1.5em]
\item \texttt{x\_halfspace\_refs::Vector\{(id,sign)\}}
\item \texttt{mapping\_id}
\item \texttt{qk\_atom\_ids}
\end{itemize}
Materialized \texttt{C\_x,C0\_x} and \texttt{C\_qK,C0\_qK} remain available for compatibility but are now derived/cache views.

\subsection{Layer E: Edge-level geometric semantics}
\texttt{VertexEdge} now records boundary provenance directly (
\texttt{x\_halfspace\_id}, \texttt{qk\_atom\_id\_from}, \texttt{qk\_atom\_id\_to}
), making graph transitions geometrically interpretable and singular-aware.

\section{Singular Regime Handling}
Nullity-1 regimes are treated as first-class. The mapping pool keeps singular metadata and $qK$-atom generation is still provenance-complete. This keeps singular/full-dimensional interfaces comparable in one data model.

\section{API Additions}
New queries expose the factorized structure directly:
\begin{itemize}[leftmargin=1.5em]
\item \texttt{get\_x\_halfspace\_refs}
\item \texttt{get\_x\_halfspace\_atom}
\item \texttt{get\_mapping\_object}
\item \texttt{get\_qk\_atom\_ids}
\item \texttt{get\_qk\_atoms}
\item \texttt{get\_qk\_atom}
\end{itemize}

\section{Why This Supports Future Rank-1 Traversal}
The key dependency is now explicit: a regime edge updates one dominance choice, therefore one boundary atom transition paired with one mapping transition. Since mappings are pooled and edge semantics are explicit, introducing Sherman--Morrison style updates can be attached to mapping transitions instead of rebuilding all constraints.

\section{Remaining Work}
A future pass should add an explicit neighbor-to-neighbor incremental mapping updater and canonical equivalence keys that deduplicate numerically equal mappings from different permutations when mathematically justified.

\end{document}
